\begin{abstract}
Self-improving agents often use task reward to decide how an experience should be written into persistent memory. This creates a direct path from reward error to future behavior. We study that path in the released ReasoningBank implementation used by a recent agent-fragility study. The code sends a trajectory with score one to a success-reflection writer. Score zero sends the identical trajectory to a failure-reflection writer. We intervene on this single reward bit while holding the released trajectory fixed. Four complete paired interventions all produce different memory content. The memory-item title set changes in every complete pair. Mean token-set Jaccard distance between paired memories is 0.735, with values from 0.691 to 0.770. The effect appears for trajectories that originally succeeded and for trajectories that originally failed. These results establish reward-conditioned memory writing as a concrete state channel: a reward-label perturbation survives the episode boundary as a different persistent memory. We then define a downstream variance experiment that injects paired memories into matched future tasks under fixed benchmark state. The protocol measures whether reward-induced memory divergence produces action divergence and success variance across repeated runs. This mechanism connects two empirical findings that have previously been studied separately: run variance in self-improving agents and reward inflation under persistent memory reuse. Reward quality is therefore part of the state-estimation problem for self-improving systems. A noisy reward signal can become a durable memory intervention before any parameter update occurs.
\end{abstract}
